\documentclass[11pt,a4paper]{article}

\usepackage{graphicx}
\usepackage{xcolor}
\usepackage[left=2.5cm, right=2.5cm, top=2.5cm, bottom=2.5cm]{geometry}
\usepackage{amsmath,amssymb,amsthm}
\usepackage{fontspec}
\setmainfont{Liberation Serif}
\setsansfont{Liberation Sans}
\setmonofont{Liberation Mono}[Scale=0.86]
\usepackage[russian]{babel}
\usepackage{booktabs}
\usepackage{tabularx}
\newcolumntype{Y}{>{\raggedright\arraybackslash}X}
\usepackage{colortbl}
\usepackage{adjustbox}
\usepackage{enumitem}
\usepackage{tcolorbox}
\definecolor{linkblue}{HTML}{1e3a5f}
\definecolor{accgreen}{HTML}{0e7c66}
\definecolor{warnamber}{HTML}{b45309}
\definecolor{rowlight}{HTML}{f0f5fa}
\usepackage[
    colorlinks=true,
    linkcolor=linkblue,
    citecolor=darkgray,
    urlcolor=linkblue,
    bookmarks=true,
    bookmarksnumbered=true,
    unicode=true
]{hyperref}

\widowpenalty=10000
\clubpenalty=10000
\emergencystretch=1.8em
\tolerance=2500

\newtcolorbox{honestybox}[1][]{colback=warnamber!6, colframe=warnamber!80,
  title={\textbf{Честность:} #1}, fonttitle=\small, boxrule=0.7pt, arc=2pt}
\newtcolorbox{findbox}[1][]{colback=accgreen!7, colframe=accgreen!80,
  title={\textbf{Находка:} #1}, fonttitle=\small, boxrule=0.7pt, arc=2pt}

\hypersetup{
    pdftitle={Открытые задачи фреймворка Чоптюка--Исаева: численные эксперименты},
    pdfauthor={Super Z (Z.ai)},
    pdfsubject={Мост Ш.3(iii), DSI-4, башня Гурвица, QNM, QCD},
    pdfkeywords={Choptuik, Choptyuk, Klein quartic, Leavitt, DSI, Hurwitz, QNM, GUE}
}

\begin{document}

\tableofcontents
\newpage

% ==========================================================================
\section{Введение и постановка задач}\label{sec:intro}
% ==========================================================================

Настоящий отчёт фиксирует первый цикл вычислительной работы над \emph{открытыми
задачами} фреймворка Чоптюка--Исаева, изложенного в монографиях репозитория
\texttt{choptuik\_ac\_bc}. Напомним цепочку, замыкающую задачу Чоптюка: спектральная
щель на квартике Клейна описывается объединённой формулой
\begin{equation}\label{eq:ch}
\Delta_{\mathrm{Ch}} \;=\; \lambda_1(D^2_{\sigma_0}) + \frac{\delta_C^2}{2} - \frac{\delta_C^5}{22},
\qquad \lambda_1(D^2_{\sigma_0}) = 3.338,\quad \delta_C = \frac{\pi}{7},
\end{equation}
воспроизводящей наблюдаемое значение $\Delta = 3.443$ с отклонением $0.15\%$;
константа фреймворка $b_{\mathrm{Ch}} = 1-\cos(2\pi/7) \approx 0.3765$ совпадает
с критическим показателем коллапса Чоптюка в пределах $0.13\%$. Все числа ниже
получены независимым прогоном пяти экспериментов
(\texttt{choptuik\_experiments/}), воспроизводимых одной командой
\texttt{python3 run\_all.py}; зёрна генераторов зафиксированы, результаты
детерминированы.

Отчёт закрывает два направления, согласованных в постановке работы. Во-первых,
\emph{открытые задачи}: мост Ш.3(iii) --- единственное аналитическое звено
верификационного приложения <<Аудит и перенос>>, остающееся постановкой, а не
результатом (эксперимент~\ref{sec:sh3}); и наблюдение DSI-4 $c_{K3} = 27/672$,
для которого мы усиливаем статус арифметикой подходящих дробей
(эксперимент~\ref{sec:dsi4}). Во-вторых, \emph{расширения фреймворка}: башня
поверхностей Гурвица с аудитом данных репозитория (эксперимент~\ref{sec:hurwitz}),
расширенный QNM-каталог LIGO/Virgo с лестницей обертонов
(эксперимент~\ref{sec:qnm}) и QCD-мост с тестом $N$-скейлинга и расширенным
сканированием $\kappa_T$ (эксперимент~\ref{sec:qcd}).

По ходу работы обнаружены три новых факта, помеченные как находки D1, E1 и
«$9/224$ --- подходящая дробь»; каждый из них сопровождается явным указанием
статуса (наблюдение / проверка / гипотеза) в духе стиля честности, принятого
в \texttt{audit\_transfer/}. Все таблицы и рисунки генерируются кодом; при
пересборке отчёта они обновляются автоматически.

% ==========================================================================
\section{Эксперимент 1. Мост Ш.3(iii): лестница препятствий}\label{sec:sh3}
% ==========================================================================

\subsection{Постановка}

Лемма Ш.3 в приложении <<Аудит и перенос>> сформулирована из трёх частей: (i)
острая теорема следа $F \ge 5n/7$ --- доказана для всех $n$; (ii) конструкция
$X_1X_1^* = X_2X_2^* = (4/7)I$ --- точный случай равенства; (iii) \emph{мост}:
софичность группы $G = U(L_{F_2}(1,2))$ индуцировала бы пары $X_1, X_2 \in
M_{n_k}(\mathbb{C})$ с дефектом на душу
\begin{equation}\label{eq:defect}
e_k \;=\; \frac{F(X_1,X_2)}{n_k} \;\le\; \eta(\varepsilon_k) \;\to\; 0,
\end{equation}
где
\begin{equation}\label{eq:Fdef}
F(X_1,X_2) = \|X_1^*X_1 - I\|^2 + \|X_2^*X_2 - I\|^2 + \|X_1^*X_2\|^2 + \|X_1X_1^* + X_2X_2^* - I\|^2
\end{equation}
--- кунцевский дефект (норма Фробениуса),
с универсальной амлификацией $\eta(\varepsilon) = C\varepsilon^{1/2}$. Поскольку
по (i) $e \ge 5/7$ в каждой $M_n$, существование такого моста противоречило бы
теореме --- в этом и состоит планируемое доказательство несофичности $G$ (цепочка
Ш1--Ш4). Часть (iii) остаётся \emph{точной постановкой}. Наш эксперимент
отвечает на практический вопрос: видят ли \emph{естественные} переносы
перестановочных данных в матрицы убывающий дефект? Мы строим «лестницу
препятствий» из пяти семейств и измеряем $e$ в каждом.

\subsection{Методы и результаты}

\textbf{(A) Теорема при больших $n$.} Тождество $F = G(P,Q)$ с $P = X_1X_1^*$,
$Q = X_2X_2^*$ проверено на $300$ случайных парах в каждой размерности
$n = 7,\dots,12$ (максимальная относительная ошибка на уровне машинной);
граница Хаара $G((\mathrm{tr}P/n)I, (\mathrm{tr}Q/n)I) \le G(P,Q)$ не нарушена
ни на одной паре. L-BFGS-минимизация с $15$ рестартами даёт
$|F_{\min} - 5n/7| \le 5.9\cdot10^{-8}$ при $n = 12$ и машинный ноль при
$n \le 10$ --- острая теорема подтверждена за пределами $n \le 6$, проверенных
в \texttt{audit\_transfer}.

\textbf{(B--D) Перестановочные семейства.} Для матриц перестановок $P_1, P_2$
дефект считается точно: $F = 2n$, т.е. $e = 2$ (подтверждено на $50$ парах для
$n = 16, 64, 128$ с ошибкой не более $5.7\cdot10^{-14}$). Масштабированное
семейство $X_i = cP_i$ даёт
\begin{equation}\label{eq:scaled}
\frac{F}{n} = 7t^2 - 8t + 3, \qquad t = c^2,
\end{equation}
с минимумом ровно $5/7$ при $t = 4/7$ --- случай равенства теоремы (i).
Двухпараметрическое семейство $X_1 = aP_1$, $X_2 = bP_2$ воспроизводит ту же
скалярную задачу $h(u,v)$ с седловой точкой $u = v = 4/7$; численно получено
$F/n = 0.7142857143$. Это замкнутая проверка: перестановки унитарны, и их
масштабирование исчерпывается теоремой (i).

\textbf{(E--F) Дерево Кунца.} Усечённые сдвиги $S_1, S_2$ двоичного дерева
глубины $k$ ($n = 2^{k+1}-1$) дают \emph{точно}
\begin{equation}\label{eq:tree}
F(S_1, S_2) = n + 2, \qquad e = 1 + \frac{2}{n} \;\to\; 1,
\end{equation}
поскольку диапазоны $S_1S_1^*$ и $S_2S_2^*$ разбивают подпространство глубины
$\ge 1$ без перекрытия ($S_1S_1^* + S_2S_2^* = P(\text{depth} \ge 1)$), так что
последний член (\ref{eq:defect}) равен $\|-P_{\mathrm{root}}\|^2 = 1$, а первые
два --- по $2^k$ (граница дерева). Взвешенное семейство $X_i = c_iS_i$ даёт
точно $F = (n-1)\big[(c_1^2-1)^2 + (c_2^2-1)^2\big] + n + 2$ (проверено до
$4.5\cdot10^{-13}$); минимум достигается при $c = (1,1)$ --- \emph{никакое
взвешивание не улучшает обрезку}. Дерево приближает отношения Ливитта
качественно лучше перестановок ($e \to 1$ против $e = 2$), но пола $5/7$ не
достигает и к нулю не идёт.

\textbf{(G) Софические данные.} Моделируя условие C1 леммы Ш.3, мы строим пары
подстановок с долей неподвижных точек $\varepsilon$ и применяем естественный
перенос $X_i = \sqrt{4/7}\,P_i$. Дефект \emph{растёт} с $\varepsilon$:
$e(0) = 5/7$, $e(0.01) = 0.745$, $e(0.05) = 0.843$, $e(0.20) = 1.140$. Закон
$\eta(\varepsilon) = C\varepsilon^{1/2}$ в этом семействе не наблюдается ---
наоборот, минимум дефекта достигается при $\varepsilon = 0$.

\begin{table}[htbp]
\centering\small
\caption{Лестница препятствий моста Ш.3(iii): дефект на душу по семействам.}
\label{tab:ladder}
\begin{tabularx}{\textwidth}{lccY}
\toprule
Семейство & $e = F/n$ & Статус & Комментарий \\
\midrule
Масштабированные унитарии $c\,U_i$ & $5/7 = 0.7143$ & теорема (равенство) & не приближает отношения Ливитта \\
Обрезка дерева Кунца & $1 + 2/n \to 1$ & точно, (\ref{eq:tree}) & лучшее «честное» семейство \\
Перестановочные модели & $2$ & точно & диапазоны перекрываются дважды \\
Софические данные (перенос) & $\ge 5/7$, растёт с $\varepsilon$ & численно, $n=256$ & $\sqrt{\varepsilon}$-закона нет \\
\bottomrule
\end{tabularx}
\end{table}

\begin{figure}[htbp]
\centering
\includegraphics[width=0.86\textwidth]{figs_ru/fig_sh3_ladder.png}
\caption{Лестница препятствий: дефект на душу $e = F/n$ для трёх семейств
против острого пола теоремы следа.}
\label{fig:ladder}
\end{figure}

\begin{figure}[htbp]
\centering
\includegraphics[width=0.86\textwidth]{figs_ru/fig_sh3_hamming.png}
\caption{Эксперимент 1G: дефект естественного переноса при софических данных с
ошибкой Хэмминга $\varepsilon$; пунктир --- гипотетический закон
$\eta(\varepsilon) = C\sqrt{\varepsilon}$.}
\label{fig:hamming}
\end{figure}

\subsection{Вывод эксперимента 1}

Все доступные конструктивные семейства платят дефект на душу не ниже $1$
(дерево) или $2$ (перестановки); пол $5/7$ достижим только тривиальным
масштабированием унитариев, которое \emph{не} приближает отношения Ливитта.
Естественные переносы не дают ни убывания дефекта, ни $\sqrt{\varepsilon}$-закона.
Это численный аргумент в пользу того, что мост Ш.3(iii) --- глубоко
нетривиальная аналитическая программа, а её итоговое отсутствие и есть
доказательство несофичности $U(L_{F_2}(1,2))$ по цепочке Ш1--Ш4.

\begin{honestybox}[статус результатов]
Доказано здесь точно: $F = n+2$ для обрезки дерева (с доказательством по
членам), $F = 2n$ для перестановок, формула (\ref{eq:scaled}) и минимум
$(u,v) = (4/7,4/7)$. Численно: подтверждение теоремы при $n \le 12$ и рост
$e(\varepsilon)$. \emph{Не} утверждается: несуществование моста (это
аналитическое утверждение); эксперимент лишь показывает, что стандартные
конструкции его не дают.
\end{honestybox}

% ==========================================================================
\section{Эксперимент 2. DSI-4: наблюдение $c_{K3} = 27/672$}\label{sec:dsi4}
% ==========================================================================

\subsection{Постановка}

Амплитуда лог-периодической модуляции критического каскада измерена в
репозитории: $c_{K3} = 0.04017757639214903$. На ведущем порядке она выводится
из DSI-ренормализации как $b_{\mathrm{Ch}}(22) = 1 - \cos(2\pi/22) = 0.0405070$;
остаток $+0.82\%$ объясняется систематикой окна измерения. Наблюдение DSI-4
фиксирует, что целочисленная комбинация
\begin{equation}\label{eq:dsi4}
c_{K3} \;=\; \frac{3^3}{2^2\cdot|\mathrm{PSL}(2,7)|} \;=\; \frac{27}{672} \;=\; \frac{9}{224} \;=\; 0.040178571\ldots
\end{equation}
воспроизводит измеренное значение с отклонением $+0.0025\%$ --- на три порядка
точнее структурных альтернатив. Статус наблюдения: не доказательство. Наша
задача --- усилить или ослабить этот статус независимыми проверками.

\subsection{Результаты}

\textbf{Цепная дробь.} Разложение измеренного $c_{K3}$ в цепную дробь даёт
$[0; 24, 1, 8, 19, 1, \dots]$, откуда подходящие дроби
$1/24$, $1/25$, \emph{$9/224$}, $172/4281$, \dots. Таким образом $9/224 = 27/672$
--- \emph{подходящая дробь измеренного значения}: она не выбрана из
бесконечного пула рациональных приближений, а arises на третьем шаге алгоритма
цепных дробей. Гауссова мера качества $|c - p/q|\cdot q^2 = 0.0499 \ll 0.5$
--- приближение исключительное. Среди всех $p/q$ с $q \le 224$ ровно одна дробь
попадает так же близко (сама $9/224$), тогда как для «случайной» цели ожидание
$\approx 0.05$ таких дробей.

\textbf{Структурные разложения.} Одно и то же число читается четырьмя
независимыми способами: $3^3/(2^2\cdot168)$ --- кубическая спинорная структура
$B^3 = -I$ над квадратичным вкладом $A^2 = -I$, нормированным на порядок группы
автоморфизмов; $9/224 = 3^2/(2^5\cdot 7)$; $(N-1)/(24N)$ при $N = 28$ ---
размерность гильбертова пространства фреймворка поверх $24$ семиугольников
тилинга $\{7,3\}$ квартики Клейна; $27/(12\cdot 56)$ поверх $56$ треугольников
того же тилинга. Каждое из разложений даёт ровно $0.040178571\ldots$

\textbf{Таблица кандидатов.} Среди тринадцати проверенных кандидатов только
$27/672$ даёт отклонение лучше $0.4\%$; бессруктурный $1/25$ отстаёт на два
порядка. Совместимость с систематикой окна: отношению
измеренное/кандидат $= 0.99998$ для $27/672$ соответствует практически
нулевой дефицит окна, тогда как $b_{\mathrm{Ch}}(22)$ требует дефицита
$0.81\%$ (внутри наблюдаемой полосы $[0.987, 0.999]$). Оба прочтения совместимы
с протоколом измерения; $27/672$ при этом в $300$ раз точнее.

\begin{table}[htbp]
\centering\small
\caption{Лучшие кандидаты амплитуды $c_{K3}$ (полная таблица --- в JSON).}
\label{tab:cands}
\begin{tabularx}{\textwidth}{lccY}
\toprule
Кандидат & Значение & Откл., \% & Статус \\
\midrule
$27/672 = 3^3/(2^2\cdot|\mathrm{PSL}(2,7)|)$ & $0.040178571$ & $+0.0025$ & наблюдение DSI-4, усилено \\
$1/25$ & $0.040000000$ & $-0.442$ & бессруктурный \\
$b_{\mathrm{Ch}}(22(1+\gamma))$ (торможённая RG) & $0.040359588$ & $+0.452$ & DSI-2 \\
$b_{\mathrm{Ch}}(22)$ & $0.040507026$ & $+0.820$ & DSI-1 \\
$1-\cos(2\pi/24)$ & $0.034074174$ & $-15.19$ & отвергнут \\
\bottomrule
\end{tabularx}
\end{table}

\begin{findbox}[$9/224$ --- подходящая дробь]
Наблюдение DSI-4 усилено: $27/672$ не просто «хорошее приближение», а элемент
последовательности подходящих дробей измеренного значения с гауссовой мерой
$0.0499$. Аналитический вывод формулы $3^3/(2^2\cdot168)$ из DSI-ренормализации
остаётся открытой задачей; теперь у неё есть точная числовая мишень.
\end{findbox}

\begin{figure}[htbp]
\centering
\includegraphics[width=0.92\textwidth]{figs_ru/fig_dsi4_candidates.png}
\caption{Структурные кандидаты амплитуды $c_{K3}$: модуль отклонения от
измеренного значения (лог-шкала).}
\label{fig:cands}
\end{figure}

% ==========================================================================
\section{Эксперимент 3. Башня поверхностей Гурвица}\label{sec:hurwitz}
% ==========================================================================

\subsection{Аудит данных репозитория: находка D1}

Монография (Часть 7 \texttt{verify\_enhanced}) использует башню
$\{$Klein $g{=}3$, $n{=}7$; Macbeath $g{=}7$, $n{=}9$; Hurwitz\_3 $g{=}14$,
$n{=}11\}$. Мы проверили каждую строку двумя обязательными соотношениями для
поверхности Гурвица: $|\mathrm{Aut}| = 84(g-1)$ и $n \mid |\mathrm{Aut}|$
(элемент порядка $n$). Первые две строки согласованы: $\mathrm{PSL}(2,7)$,
$168 = 84\cdot 2$, $7\mid 168$; $\mathrm{PSL}(2,8)$, $504 = 84\cdot 6$,
$9\mid 504$. Третья строка \textbf{несогласована}: $84(14-1) = 1092 =
2^2\cdot 3\cdot 7\cdot 13$ не делится на $11$, поэтому группа порядка $1092$ не
имеет элементов порядка $11$, и сигнатура $(2,3,11)$ для рода $14$ невозможна.

\begin{findbox}[D1: несогласованность башни репо]
Для $g = 14$ реальная поверхность Гурвица --- фактор по $\mathrm{PSL}(2,13)$,
$|\mathrm{PSL}(2,13)| = 1092$, сигнатура $(2,3,13)$. В корректной башне
$\delta = \pi/13$, а не $\pi/11$; все вычисления фреймворка для этой точки
нужно пересчитать. Мы считаем это редакторской ошибкой данных (не математики
поправок).
\end{findbox}

\subsection{Корректная башня и расчёты}

Мы строим корректную башню из четырёх младших поверхностей Гурвица, включая
\emph{первый триплет} рода $17$ ($|G| = 1344 = 84\cdot 16$, сигнатура
$(2,3,8)$) --- ближайшую после Клейна точку, где башня ветвится. Все проверки
$84(g-1) = |\mathrm{Aut}|$, $n \mid |\mathrm{Aut}|$ и гиперболичности
$1/2 + 1/3 + 1/n < 1$ проходят. Для каждой поверхности вычислены фазовые
величины фреймворка $\delta = \pi/n$, $\gamma = \delta^4/b_2$,
$\delta_{\mathrm{eff}} = \delta^5/b_2$ при гипотетическом паттерне
$b_2(k) = 2b_2(k-1) + 2 = 3\cdot 2^{k+3} - 2$, воспроизводящем три значения
репо $(22, 46, 94)$ и продолжающем их на четвёртую точку ($190$).

\begin{table}[htbp]
\centering\small
\caption{Корректная башня Гурвица и величины фреймворка. $b_2(k)$ --- гипотеза
паттерна (для Клейна топологична: $b_2(K3) = 22$).}
\label{tab:tower}
\begin{tabular}{llcclll}
\toprule
Поверхность & Группа & $g$ & $n$ & $|\mathrm{Aut}|$ & $\gamma$ & $\delta_{\mathrm{eff}}$ \\
\midrule
Клейн & PSL(2,7) & 3 & 7 & 168 & $1.844\cdot10^{-3}$ & $8.276\cdot10^{-4}$ \\
Макбит & PSL(2,8) & 7 & 9 & 504 & $3.228\cdot10^{-4}$ & $1.127\cdot10^{-4}$ \\
PSL(2,13) & PSL(2,13) & 14 & 13 & 1092 & $3.628\cdot10^{-5}$ & $8.768\cdot10^{-6}$ \\
Триплет (первый) & $(2,3,8)$, $|G|=1344$ & 17 & 8 & 1344 & $1.252\cdot10^{-4}$ & $4.915\cdot10^{-5}$ \\
\bottomrule
\end{tabular}
\end{table}

Давление $\delta_{\mathrm{eff}}$ вдоль PSL-ветви составляет $\times 7.3$ (Клейн
$\to$ Макбит) и $\times 12.9$ (Макбит $\to$ PSL(2,13)); ветвление на триплете
возвращает часть амплитуды ($n$ падает с $13$ до $8$). Универсальная кривая
$\Delta_{\mathrm{Ch}}(\delta) = 3.338 + \delta^2/2 - \delta^5/22$ монотонно
растёт при $\delta \in (0, 1.64)$ (производная $\delta(1 - 5\delta^3/22) > 0$),
поэтому «минимум вблизи $\pi/7$» в монографии корректно читать как точку
пересечения кривой с наблюдаемым значением: $\delta_* = 0.4603$ против
$\pi/7 = 0.4488$ (расхождение $0.026$). Точки $\pi/n$ корректной башни дают
$\Delta_{\mathrm{Ch}}(\pi/13) = 3.3672$ --- предсказание фреймворка для род-14
спектральной величины, которое может быть сверено, когда для этой поверхности
будет известно $\lambda_1$.

\begin{honestybox}[гипотетичность $b_2(k)$]
Для не-Клейновых поверхностей башни значение $b_2$ не выведено из топологии:
топологичен только $b_2 = 22 = b_2(K3)$ для Клейна. Паттерн $b_2(k) =
3\cdot 2^{k+3}-2$ --- экстраполяция трёх точек репо; все числа таблицы
\ref{tab:tower} следует читать как \emph{структурную гипотезу} фреймворка.
Формальная экстраполяция дальше четвёртой точки ($n = 10, 12, 14$ при
$b_2 = 382, 766, 1534$) не привязана к реальным поверхностям и приведена только
как продолжение паттерна.
\end{honestybox}

\begin{figure}[htbp]
\centering
\includegraphics[width=0.88\textwidth]{figs_ru/fig_hurwitz_tower.png}
\caption{Корректная башня Гурвица: торможение $\gamma$ и эффективная фаза
$\delta_{\mathrm{eff}}$ по поверхностям (лог-шкала).}
\label{fig:tower}
\end{figure}

% ==========================================================================
\section{Эксперимент 4. QNM: каталог LIGO/Virgo и обертоны}\label{sec:qnm}
% ==========================================================================

\subsection{Постановка и метод}

Фреймворк v2.0 предсказывает мультипликативную поправку к частотам
квазинормальных мод чёрных дыр:
\begin{equation}\label{eq:qnm}
f_{\mathrm{corr}} = f\left(1 - \frac{\delta_{\mathrm{eff}}}{\pi^2}\right)
= 0.999916143\,f, \qquad
\frac{\Delta f}{f} = -8.3857\cdot 10^{-5}.
\end{equation}
Мы расширяем каталог репозитория с четырёх до десяти событий. Для четырёх
«якорных» событий используются измеренные частоты ringdown из репозитория; для
шести новых события частота фундаментальной моды $l = m = 2$ вычисляется по
подгонке Берти--Кардозо--Уилла по финальным массе и спину из публикаций
LVK: $M\omega_R = 1.5251 - 1.1568(1-a)^{0.1292}$. Для якорных событий обе
оценки приведены рядом --- расхождение $8$--$13\%$ отражает точность подгонки и
метода экстракции ringdown, а не ошибку данных.

\begin{table}[htbp]
\centering\small
\caption{Каталог из десяти событий: предсказанные сдвиги $\Delta f$.
Для якорных событий ($\ast$) $f$ измерена (репо), для остальных --- подгонка
Берти.}
\label{tab:qnm}
\begin{tabular}{lrrrrr}
\toprule
Событие & $M_f$, $M_\odot$ & $a_f$ & $f$, Гц & $\Delta f$, Гц & $f_{\mathrm{corr}}$, Гц \\
\midrule
GW150914$^*$ & 62.3 & 0.67 & 251.0 & $-0.0210$ & 250.979 \\
GW151226 & 20.5 & 0.74 & 871.8 & $-0.0731$ & 871.710 \\
GW170104$^*$ & 48.7 & 0.65 & 314.0 & $-0.0263$ & 313.974 \\
GW170608 & 17.8 & 0.69 & 963.5 & $-0.0808$ & 963.419 \\
GW170729 & 59.1 & 0.81 & 323.5 & $-0.0271$ & 323.473 \\
GW170809 & 35.4 & 0.70 & 488.3 & $-0.0409$ & 488.259 \\
GW170814$^*$ & 53.4 & 0.70 & 286.0 & $-0.0240$ & 285.976 \\
GW170818 & 35.6 & 0.67 & 474.4 & $-0.0398$ & 474.360 \\
GW190412 & 40.3 & 0.87 & 510.2 & $-0.0428$ & 510.157 \\
GW190521$^*$ & 142.0 & 0.72 & 110.0 & $-0.0092$ & 109.991 \\
\bottomrule
\end{tabular}
\end{table}

\subsection{Обнаружимость и обертоны}

Поскольку поправка (\ref{eq:qnm}) мультипликативна, отношение $|\Delta f|/\sigma$
определяется только относительной прецизионностью детектора: $0.004\sigma$ при
текущих $2\%$, $0.084\sigma$ для Einstein Telescope ($0.1\%$) и $0.84\sigma$
для Cosmic Explorer ($0.01\%$). Порог $1\sigma$ достигается при прецизионности
$0.0084\%$; объединение $N$ независимых событий даёт выигрыш $\sqrt{N}$, так
что каталог из десяти событий уже на CE-классе детекторов в сумме превышает
порог. Это уточняет формулу репозитория <<near-detectable>> и делает предсказание
фальсифицируемым на горизонте детекторов третьего поколения.

Лестница обертонов для GW150914 построена в пределе Шварцшильда по эталонным
значениям $M\omega_n$ (Берти и др., 2006): $f_0..f_4 = 251,\ 394,\ 537,\ 681,\
825$ Гц, времена затухания $\tau_0..\tau_4 = 3.45,\ 1.57,\ 1.02,\ 0.76,\
0.60$ мс. Мультипликативность поправки даёт фальсифицируемое предсказание:
$\Delta f_n / f_n = \mathrm{const} = 8.3857\cdot10^{-5}$ \emph{для всех}
обертонов --- относительные интервалы $f_n/f_0$ не меняются. Спектроскопия
обертонов, обнаруживающая непропорциональные сдвиги, фальсифицировала бы
фреймворк; пропорциональные сдвиги с указанной константой --- подтвердили бы.

\begin{honestybox}[статус QNM-чисел]
Для не-якорных событий $f$ --- GR-предсказание по $(M_f, a_f)$, а не измерение;
поправка применяется к GR-значению. Обертоны --- в пределе Шварцшильда: при
$a = 0.67$ частоты сдвинуты на $O(1-a) \sim 30\%$; лестница используется как
схема масштабирования, а не как точная таблица Керра.
\end{honestybox}

\begin{figure}[htbp]
\centering
\includegraphics[width=0.88\textwidth]{figs_ru/fig_qnm_overtones.png}
\caption{Лестница обертонов GW150914 ($l = m = 2$, предел Шварцшильда,
масштаб к измеренной $f_0 = 251$ Гц): частоты $f_n$ и сдвиги $\Delta f_n$.}
\label{fig:overtones}
\end{figure}

% ==========================================================================
\section{Эксперимент 5. QCD-мост: $N$-скейлинг и $\kappa_T$-sweep}\label{sec:qcd}
% ==========================================================================

\subsection{Находка E1: след оператора $O_\chi$}

Work-формула фреймворка $\bar\theta_{\mathrm{eff}} = \delta_C\, N\langle\lambda\rangle
\mathcal{S}_{\mathrm{GUE}}$ опирается на то, что в GUE-классе спектр симметричен
и $\langle\lambda\rangle = 0$ в континууме. Мы построили \emph{башню}
$O_\chi(N = 28k)$ из $k$ блоков $(Q_{K3} \oplus M_F)$ с полной матрицей
$T$-нарушения $V_T$ и обнаружили: тривиальный след ядра
$\mathrm{tr}(Q_{K3}\oplus M_F) = 32$ (диагонали матриц Картана $E_8$ равны $2$),
поэтому сырое среднее $\langle\lambda\rangle = 32/28 = 1.143$ --- \emph{константа}
по всей башне, и аргумент симметрии к сырому оператору неприменим. Поправка
одна строка и физически обязательна: $\bar\theta$ сопряжён с \emph{флуктуацией}
топологического заряда, а не с его постоянной составляющей, поэтому оператор
следует центрировать: $O_\chi \to O_\chi - (\mathrm{tr}O_\chi/N)I$. После
центрирования флуктуация выборочного среднего убывает по степенному закону с
измеренным наклоном $-0.99 \pm$ --- т.е. $|\langle\lambda\rangle| \sim 1/N$,
даже быстрее заявленного $1/\sqrt{N}$.

\begin{findbox}[E1: центрирование $O_\chi$]
В конструкции $O_\chi$ репозитория след $32$ не убран. Для work-формулы
необходимо центрирование; после него $\langle\lambda\rangle \to 0$ с наклоном
$-0.99$ по $\log N$ (башня $N = 28..448$, $M = 48$ реализаций на точку). Это
редакторская правка конструкции, усиливающая вывод фреймворка.
\end{findbox}

\subsection{$\kappa_T$-sweep: переход в GUE-класс}

Расширенное сканирование $\kappa_T \in [0, 16]$ на $N = 28$ ($200$ реализаций
$V_T$ на точку, эталонные ансамбли GUE/GOE/Poisson по $3000$ выборок)
воспроизводит и уточняет картину репозитория. При $\kappa_T = 0$ спектр
вырожден (целочисленный блочный оператор); с ростом связи статистика
расстояний монотонно переходит от пуассоновской к GUE: $\bar r$ растёт от
$0.045$ до $\approx 0.36$ (эталон GUE $0.387$, Poisson $1/3$). Байес-фактор
$\mathrm{BF}(\mathrm{GUE}/\mathrm{Poi})$ достигает «убедительного» уровня
($\ge 100$) при $\kappa_T \approx 4$ и максимален при $\kappa_T = 8.45$
($\mathrm{BF} \approx 1165$), согласуясь с решёточной оценкой
$\kappa_T^{(\mathrm{QCD})} > 2.62$ (95\% CL), при которой репозиторий
интерполирует $\mathrm{BF} \ge 99$. В нашей независимой сетке при
$\kappa_T = 2.62$ получено $\mathrm{BF} = 50$ (strong); качественный вывод
(«при физическом $\kappa_T$ оператор в GUE-классе») воспроизводится.

\begin{table}[htbp]
\centering\small
\caption{$\kappa_T$-sweep на $N = 28$ после центрирования (находка E1):
среднее фолдинговое отношение и байес-фактор.}
\label{tab:kappa}
\begin{tabular}{rcccl}
\toprule
$\kappa_T$ & $\bar r$ & $\sigma_r$ & BF(GUE/Poi) & Вердикт \\
\midrule
0.00 & 0.045 & 0.000 & $0.0$ & вырожден \\
0.50 & 0.243 & 0.020 & $0.0$ & Poisson-класс \\
1.00 & 0.308 & 0.022 & $0.8$ & переход \\
1.50 & 0.333 & 0.023 & $15.5$ & STRONG GUE \\
2.62 & 0.344 & 0.022 & $50.3$ & STRONG GUE \\
4.00 & 0.353 & 0.023 & $193.3$ & DECISIVE GUE \\
8.45 & 0.362 & 0.024 & $1164.9$ & DECISIVE GUE \\
16.00 & 0.358 & 0.024 & $501.8$ & DECISIVE GUE \\
\bottomrule
\end{tabular}
\end{table}

\begin{honestybox}[модель $V_T$ и сетка BF]
$V_T$ остаётся модельной матрицей $T$-нарушения (гауссов эрмитов блок с
масштабом $b_C$) --- конструктивной гипотезой фреймворка. Числа BF зависят от
сетки и ядер KDE; сравнение с интерполяциями репозитория ($99$ при $2.62$,
$510$ при $8.45$) --- качественное, не численное.
\end{honestybox}

\begin{figure}[htbp]
\centering
\includegraphics[width=0.86\textwidth]{figs_ru/fig_qcd_nscaling.png}
\caption{$N$-скейлинг центрированной флуктуации $\langle\lambda\rangle$ на
башне $O_\chi$: измеренный наклон $-0.99$ против эталонов $1/\sqrt{N}$ и
$1/N$.}
\label{fig:nscaling}
\end{figure}

\begin{figure}[htbp]
\centering
\includegraphics[width=0.86\textwidth]{figs_ru/fig_qcd_kappa_sweep.png}
\caption{$\kappa_T$-sweep: байес-фактор BF(GUE/Poisson); вертикаль ---
решёточный порог $\kappa_T > 2.62$ (95\% CL).}
\label{fig:kappa}
\end{figure}

% ==========================================================================
\section{Честность и ограничения}\label{sec:honesty}
% ==========================================================================

Доказано аналитически (внутри эксперимента): точная формула $F = n + 2$ для
обрезки дерева Кунца с разбиением диапазонов; точная формула
$F = (n-1)[(c_1^2-1)^2 + (c_2^2-1)^2] + n + 2$ для взвешенного дерева и её
минимум при $c = (1,1)$; формула $F/n = 7t^2 - 8t + 3$ для масштабированных
перестановок с минимумом $5/7$ при $t = 4/7$; цепочка равенств
(\ref{eq:dsi4}). Проверено численно с указанной точностью: острая теорема следа
при $n \le 12$ (L-BFGS, зазор $\le 5.9\cdot10^{-8}$); исключительность $9/224$
(гауссова мера $0.0499$; одна дробь при $q \le 224$ против ожидания $0.05$);
рост $e(\varepsilon)$ при софических данных; степенной закон центрированной
флуктуации $\langle\lambda\rangle$ (наклон $-0.99$).

Не утверждается и остаётся открытым: несуществование моста Ш.3(iii) (наше
результат --- только недоказуемость его стандартными конструкциями);
аналитический вывод $27/672$ из DSI-ренормализации (у задачи появилась точная
мишень --- подходящая дробь); вывод $b_2(k)$ башни Гурвица из топологии;
модель $V_T$ в QCD-экспериментах. Находки D1 (несогласованность $n = 11$ в
башне репо) и E1 (след $32$ в $O_\chi$) сформулированы как редакторские правки
данных и конструкции, не как опровержение математики поправок.

% ==========================================================================
\section{Выводы и дальнейшие шаги}\label{sec:concl}
% ==========================================================================

Выполнен первый полный цикл работы над открытыми задачами фреймворка. По мосту
Ш.3(iii) построена лестница препятствий: между лучшим «честным» семейством
(дерево Кунца, $e \to 1$) и острым полом ($5/7$) остаётся зазор, который не
закрывается ни перестановками, ни взвешиванием, ни софическими данными с
естественным переносом --- тем самым аналитическая программа (iii) получает
численный контур целевой конструкции. По DSI-4 наблюдение $27/672$ усилено
статусом подходящей дроби с гауссовой мерой $0.0499$ и четырьмя независимыми
структурными разложениями. По башне Гурвица найдена и исправлена
несогласованность данных репозитория (D1) и построена корректная четырёхточечная
башня с предсказанием $\Delta_{\mathrm{Ch}}(\pi/13) = 3.3672$. QNM-каталог
расширен до десяти событий с явным горизонтом фальсификации (прецизионность
$0.0084\%$, детекторы CE-класса) и фальсифицируемым обертоновым предсказанием
$\Delta f_n/f_n = \mathrm{const}$. В QCD-эксперименте найдена необходимая
правка центрирования $O_\chi$ (E1), после которой $N$-скейлинг подтверждён с
наклоном $-0.99$, а $\kappa_T$-sweep воспроизводит GUE-переход с
$\mathrm{BF} \ge 50$ уже на решёточном пороге.

Дальнейшие шаги вытекают из находок. Во-первых, внести правки D1 и E1 в
репозиторий (таблица Hurwitz\_3 и конструкция $O_\chi$) и пересобрать
верификацию. Во-вторых, атаковать аналитическую мишень DSI-4: объяснить,
почему DSI-ренормализация с $\lambda = 22$ выделяет именно подходящую дробь
$9/224$; естественный путь --- исследовать, фиксирует ли окно измерения
знаменатель $224 = 4\cdot 56$ через комбинаторику тилинга $\{7,3\}$. В-третьих,
замкнуть зазор $[5/7,\,1]$ по мосту Ш.3(iii): доказать нижнюю оценку $e \ge 1$
для класса «изометрических на носителе» конструкций, к которым относится дерево,
--- это отделило бы обрезку дерева от масштабированных унитариев и превратило
бы лестницу препятствий в теорему. Наконец, QNM-направление выигрывает от
пересчёта обертонов в керровском пределе ($a = 0.67$) по точным таблицам, чтобы
заменить схему масштабирования предсказанием полной точности.

% ==========================================================================
\appendix
\section{Воспроизводимость}\label{sec:repro}
% ==========================================================================

Все эксперименты выполнены кодом каталога experiments (Python 3.12,
NumPy/SciPy/Matplotlib); оркестратор \texttt{run\_all.py} запускает пять
скриптов одной командой, полное время прогона --- порядка пяти минут. Зёрна
генераторов зафиксированы: 42 --- эталонные ансамбли QCD; 5000 + 97r + n ---
рестарты L-BFGS; 9000 + 131r + k --- дерево; 11000 + \ldots{} --- софические
данные. Два прогона подряд дают побитово совпадающие файлы JSON. Результаты
лежат в каталоге results (файлы exp*.json), а фигуры --- в каталогах
fig\_ru и fig\_en, по одному комплекту на каждый язык отчёта.
Проверки целостности: тестовый набор репозитория (30 тестов), верификация
\texttt{audit\_transfer} и полная верификация \texttt{verify\_enhanced}
пройдены перед началом экспериментов на том же окружении.

\end{document}
